% Generated draft: requires mathematical and source review.
\documentclass[11pt]{article}
\usepackage[margin=25mm]{geometry}
\usepackage{fontspec}
\setmainfont{Latin Modern Roman}
\usepackage{amsmath,amssymb,mathtools,unicode-math}
\setmathfont{Latin Modern Math}
\usepackage{enumitem,needspace,xurl,hyperref}
\setlength{\parindent}{0pt}
\setlength{\parskip}{5pt}
\setlength{\emergencystretch}{3em}
\newcommand{\N}{\mathbb{N}}
\newcommand{\Z}{\mathbb{Z}}
\newcommand{\Q}{\mathbb{Q}}
\newcommand{\R}{\mathbb{R}}
\newcommand{\C}{\mathbb{C}}
\newcommand{\F}{\mathbb{F}}
\newcommand{\A}{\mathbb{A}}
\newcommand{\PP}{\mathbb{P}}
\newcommand{\E}{\mathbb{E}}
\newcommand{\eps}{\varepsilon}
\newcommand{\dd}{\,\mathrm{d}}
\begin{document}
\section*{UnsolvedMath ID: 30003710}
\textbf{Source problem code:} catalogue entry
\subsection{Definitions and mathematical statement}
Let \(i\in\Z_{\geq 0}\) and let \(p\) be a prime. For \(n\geq 1\), let \(F_n\) be the free group of rank \(n\), and write \(\operatorname{Out}(F_n)=\operatorname{Aut}(F_n)/\operatorname{Inn}(F_n)\). Put \(\F_p=\Z/p\Z\). For a group \(G\) and a specified \(G\)-module \(M\), write \(H^i(G;M)\) for group cohomology.

The supplied question asks whether
\[
H^i\bigl(\operatorname{Out}(F_n);\F_p^{\,n}\bigr)
\]
is eventually periodic in \(n\), for every fixed \(i\) and \(p\). The coefficient action is not explicitly specified. A natural interpretation identifies \(\F_p^{\,n}\) with
\[
V_n=(F_n/[F_n,F_n])\otimes_{\Z}\F_p,
\]
with the action induced by automorphisms of \(F_n\); inner automorphisms act trivially on this quotient, so this gives an \(\operatorname{Out}(F_n)\)-module. This interpretation requires confirmation from the source.

Under this coefficient interpretation, and interpreting eventual periodicity as periodicity of isomorphism classes of \(\F_p\)-vector spaces, the question is whether
\[
\forall i\in\Z_{\geq0}\ \forall p\text{ prime}\ \exists N\geq1\ \exists T\geq1\ \forall n\geq N,\qquad
H^i\bigl(\operatorname{Out}(F_{n+T});V_{n+T}\bigr)
\cong_{\F_p}
H^i\bigl(\operatorname{Out}(F_n);V_n\bigr).
\]
Here \(N\) and \(T\) may depend on \(i\) and \(p\). The supplied record does not specify whether periodicity concerns these isomorphism classes, numerical dimensions, or isomorphisms compatible with additional maps or structure; no such maps or structure are supplied.

If the initial question has an affirmative answer, the record asks the same question for coefficients given by a polynomial functor or a VIC-module. For a specified coefficient system \((M_n)_{n\geq1}\) with specified \(\operatorname{Out}(F_n)\)-actions, the corresponding question would replace \(V_n\) by \(M_n\) in the displayed periodicity condition, allowing \(N\) and \(T\) also to depend on that system. The source record does not specify the polynomial functors, their source and target categories or degree conditions, or the base ring, morphism convention, variance, and coefficient construction for VIC-modules. Consequently, the intended class of coefficient systems and the precise quantifiers for this extension remain unresolved.
\subsection{Short English statement}
For each fixed cohomological degree and prime, ask whether the cohomology of \(\operatorname{Out}(F_n)\) with the indicated rank-\(n\) modular coefficients eventually repeats with a positive period as \(n\) grows. If so, ask the analogous question for polynomial-functor or VIC-module coefficients. The coefficient action, meaning of repetition, and scope of the generalized coefficients require confirmation.
\subsection{Sources}
\begin{itemize}
\item Local UnsolvedMath registry: permanent ID \(30003710\), source problem code \texttt{catalogue entry}, title \emph{Periodicity of Twisted Outer-Automorphism-Group Cohomology}; supplied as an Oberwolfach Reports Open Problems record. The local registry paths are \texttt{data/unsolvedmath/problems.json} and \texttt{catalog/problems.json}; this draft uses the supplied complete record.
\item Supplied report citation: \emph{Topology of Arrangements and Representation Stability (2019)}, Oberwolfach Reports, DOI \texttt{10.4171/owr/2018/2}, pp.\ 43--123. Supplied source URL: \url{https://doi.org/10.4171/owr/2018/2}. The record describes these reports as documenting mathematical workshops and their open-problem discussions. It gives proposed year 2018 and no named proposer.
\item Supplied literature-source label: \emph{Oberwolfach Report 2/2018}, \url{https://doi.org/10.4171/owr/2018/2}. The literature assessment dated exactly \texttt{2026-08-21} labels the exact modular eventual-periodicity question and its polynomial-functor or VIC variants unresolved, and states that rational representation-stability results do not answer this modular question. This is supplied dated curation derived from the OWR record, not an independently verified literature assessment; time-sensitive status and cited claims require independent verification.
\item The supplied registry labels the problem \texttt{open} and the statement \texttt{corrected\_verified}, explaining that mathematical content was preserved while notation, typography, or omitted definitions were repaired. These are supplied editorial labels, not mathematical verification of this draft.
\end{itemize}
\end{document}
